\documentclass[11pt,a4paper]{article}

\usepackage{amsmath,amssymb,amsthm}
\usepackage[margin=1in]{geometry}
\usepackage{hyperref}
\usepackage{booktabs}
\usepackage{url}

\newtheorem{theorem}{Theorem}[section]
\newtheorem{proposition}[theorem]{Proposition}
\newtheorem{lemma}[theorem]{Lemma}
\newtheorem{corollary}[theorem]{Corollary}
\theoremstyle{definition}
\newtheorem{definition}[theorem]{Definition}
\theoremstyle{remark}
\newtheorem{remark}[theorem]{Remark}

\newcommand{\tr}{\operatorname{tr}}
\newcommand{\Nil}{\operatorname{Nil}}
\newcommand{\rank}{\operatorname{rank}}
\newcommand{\disc}{\operatorname{disc}}
\newcommand{\kinv}{k_{\mathrm{inv}}}
\newcommand{\SLtwo}{\mathrm{SL}_2(\mathbb{C})}
\newcommand{\PSLtwo}{\mathrm{PSL}_2(\mathbb{C})}
\newcommand{\Zfive}{\mathbb{Z}/5}
\newcommand{\Q}{\mathbb{Q}}
\newcommand{\Z}{\mathbb{Z}}

\title{Arithmetic and Character-Variety Structure of the\\
Closed Hyperbolic 3-Manifold $m003(-2,3)$}

\author{Marvin Gentry\\
\small Independent Researcher, Seattle, Washington\\
\small \texttt{drmlgentry@protonmail.com} \quad ORCID: 0009-0006-4550-2663}

\date{September 2026}

\begin{document}

\maketitle

\begin{abstract}
We establish exact arithmetic and character-variety results for the closed
hyperbolic 3-manifold $M = m003(-2,3)$, obtained by Dehn filling the
one-cusped manifold $m003$ along the slope $(-2,3)$.

First, we prove the surgery law
$|H_1(m003(p,q);\Z)| = 5\,|2p+q|$
directly from the fundamental-group presentation, so that
$H_1(M) \cong \Zfive$ is a consequence of the topology and peripheral
structure of $m003$ rather than an imposed constraint.

Second, we determine the invariant trace field exactly:
$\kinv(M) \cong \Q[X]/(X^4+X^3-1)$,
a quartic field of discriminant $-283$ and signature $(2,1)$.
The proof uses exact Gr\"obner elimination from the filled-group
presentation together with certified interval arithmetic, and
contains no numerical polynomial recognition (\texttt{algdep} or PSLQ)
in any load-bearing step. This corrects an earlier claim in which
the invariant trace field of the closed filling was conflated with
the Eisenstein field $\Q(\sqrt{-3})$ of the cusped parent.

Third, a preregistered target-free short-word atlas over thirteen
Dehn fillings of $m003$ exposes four squared-trace degeneracies.
Exact primary decomposition shows that three of them are identities
of the geometric component $X_0$ of the $\SLtwo$-character variety,
while the fourth,
$\tr^2 \rho(B) = \tr^2 \rho(Abb)$,
characterizes the $(-2,3)$ filling locus among irreducible projective
characters. We give the group-theoretic mechanism:
$BaBA \cdot B \cdot (BaBA)^{-1} = Abb^{-1}$ in $\pi_1(M)$.

Finally, we revisit a previously proposed Borel-factor construction for
the lepton mixing matrix on this manifold. The construction is
reproducible, but a matched-procedure null calibration gives
$\hat p = 0.1325$ ($6625/50000$), so the earlier significance claim
of $p < 10^{-4}$ does not survive and is withdrawn. We report this
negative result explicitly and keep it logically separate from the
exact results above.
\end{abstract}

\tableofcontents

\section{Introduction}

The manifold $m003$ in the SnapPy cusped census is a one-cusped
finite-volume hyperbolic 3-manifold whose cusp cross-section is the
Eisenstein torus. Its Dehn fillings along a distinguished ray of slopes
all have first homology $\Zfive$, and the first member of that ray,
$m003(-2,3)$, is a closed hyperbolic 3-manifold of volume
$0.981368\ldots$

This paper reports what can be proved exactly about that manifold.
The results fall into three groups: a homology surgery law, an exact
determination of the invariant trace field, and an exact description
of certain trace degeneracies in the $\SLtwo$-character variety.

A fourth section describes a numerical construction relating holonomy
data on this manifold to the lepton mixing matrix. That section is
included for completeness and is explicitly downstream of the
mathematics: it reports a reproducible computation together with a
null calibration that does not support a claim of statistical
exceptionality.

\subsection{What this paper does not claim}

Several statements that appeared in earlier drafts of this work are
withdrawn here, and we list them so that no reader reconstructs them
from context.

\begin{itemize}
\item We do \emph{not} claim that $M$ is the minimum-volume closed
orientable hyperbolic 3-manifold. That manifold is the Weeks manifold,
of volume $0.942707\ldots$, with $H_1 \cong \Zfive \oplus \Zfive$
\cite{GMM2009}. The correct statement, used below, is that $M$ is of
minimal volume \emph{within the $H_1 \cong \Zfive$ stratum}.

\item We do \emph{not} claim that $\kinv(M) = \Q(\sqrt{-3})$. The
Eisenstein field belongs to the cusped parent $m003$; the closed
filling has the quartic field $K_{283}$ established in
Section~\ref{sec:itf}.

\item We do \emph{not} claim that a Lucas-prime covering-tower
condition or a golden-ratio scale parameter selects this manifold.
Both were investigated and neither survived scrutiny.

\item We do \emph{not} claim that the lepton mixing fit reported in
Section~\ref{sec:borel} is statistically exceptional. The corrected
null calibration in Section~\ref{sec:null} rejects that reading.
\end{itemize}

\section{The manifold and its homology}
\label{sec:homology}

\subsection{Presentation}

We fix throughout the two-generator presentation of the cusped
manifold returned by SnapPy,
\begin{equation}
\pi_1(m003) = \langle a, b \mid r \rangle,
\qquad r = \texttt{abAAbabbb},
\end{equation}
with peripheral words
\begin{equation}
\mu = \texttt{ABABB}, \qquad \lambda = \texttt{ABAbab} .
\end{equation}
All statements below refer to this marked presentation. Where a filled
manifold is used, we verify explicitly that the filling does not
change the generating set, and we certify separately that any
shortened filling relator returned by the software generates the same
normal closure as the literal slope word (Section~\ref{sec:bridge}).

\subsection{The surgery law}

\begin{theorem}[Surgery law]
\label{thm:surgery}
For coprime $(p,q)$ with $m003(p,q)$ defined,
\begin{equation}
\bigl| H_1\bigl(m003(p,q);\Z\bigr) \bigr| \;=\; 5\,|2p+q| .
\end{equation}
\end{theorem}

\begin{proof}
Abelianizing the presentation, the relator $r=\texttt{abAAbabbb}$ has
exponent vector $(0,5)$ in $([a],[b])$ (four $a$-letters cancelling in
pairs, five net $b$-letters), giving the single relation $5[b]=0$ and
hence $H_1(m003) \cong \Z\langle a\rangle \oplus \Zfive\langle
b\rangle$, with the $\Zfive$ torsion generated by the class of $b$, not
$a$. The peripheral curves $\mu=\texttt{ABABB}$, $\lambda=\texttt{ABAbab}$
abelianize to $[\mu]=-2[a]-3[b]$ and $[\lambda]=-[a]+[b]$, so filling
along $(p,q)$ adjoins the relation $p[\mu]+q[\lambda]=(-2p-q)[a]+
(-3p+q)[b]=0$. The resulting $2\times2$ integer presentation matrix
$\left(\begin{smallmatrix}0 & 5\\ -2p-q & -3p+q\end{smallmatrix}\right)$
has determinant $5(2p+q)$, giving $|H_1|=5|2p+q|$ directly; Smith
normal form confirms the group is cyclic of that order whenever
$|2p+q|=1$, in which case $H_1\cong\Zfive$.
\end{proof}

\begin{corollary}
\label{cor:z5ray}
$H_1(m003(p,q)) \cong \Zfive$ if and only if $|2p+q| = 1$. The ray
$2p+q = -1$ contains
\begin{equation}
(-2,3),\ (-3,5),\ (-4,7),\ (-5,9),\ \ldots
\end{equation}
\end{corollary}

\begin{remark}
The torsion order $5$ is therefore not a parameter chosen to fit
anything. It is forced by the relator, and the coefficient $2$ in
$2p+q$ is forced by the peripheral curves. Both are consequences of
the triangulation of $m003$.
\end{remark}

\subsection{Position of $M$ within the $\Zfive$ stratum}

Within the SnapPy \texttt{OrientableClosedCensus}, exactly $134$
manifolds have $H_1 \cong \Zfive$. Among these, $M = m003(-2,3)$
has minimal volume, $\mathrm{vol}(M) = 0.981368\ldots$

We emphasize the scope: this is minimality within a homology stratum,
not global volume minimality. The Weeks manifold has smaller volume
but lies outside the stratum, having $H_1 \cong \Zfive \oplus \Zfive$.

\subsection{The cusped parent}

The cusp shape of $m003$ is $\tau = \tfrac12 + \tfrac{\sqrt3}{2} i$,
the Eisenstein point, and
\begin{equation}
\kinv(m003) = \Q(\sqrt{-3}).
\end{equation}
This field belongs to the cusped manifold only. The invariant trace
field of the closed filling is computed in the next section and is
a different, larger field.

\section{The invariant trace field of $m003(-2,3)$}
\label{sec:itf}

\begin{theorem}[Invariant trace field]
\label{thm:itf}
Let $M = m003(-2,3)$ and
\begin{equation}
K_{283} \;=\; \Q[X]/(X^4+X^3-1).
\end{equation}
Then $X^4+X^3-1$ is irreducible over $\Q$, and
\begin{equation}
\boxed{\;\kinv(M) \;\cong\; K_{283}\;}
\end{equation}
a quartic field of discriminant $-283$ and signature $(2,1)$.
\end{theorem}

The proof follows the four-gate structure described below. No step
uses \texttt{algdep}, PSLQ, or any other numerical polynomial
recognition as a load-bearing ingredient.

\subsection{Gate F1: exact elimination from the presentation}

Using the rank-two chart
\begin{equation}
A = \begin{pmatrix} x & -1 \\ 1 & 0 \end{pmatrix},
\qquad
B = \begin{pmatrix} 0 & -u \\ u^{-1} & y \end{pmatrix},
\qquad u^2 + zu + 1 = 0,
\end{equation}
so that $x = \tr A$, $y = \tr B$, $z = \tr(AB)$, we impose the full
matrix relations
\begin{equation}
\rho(r) = I, \qquad \rho\bigl(\mu^{-2}\lambda^{3}\bigr) = I,
\end{equation}
together with $T - \tr\rho(\mu) = 0$, and eliminate
$x,y,z,u$ by lexicographic Gr\"obner basis. This produces a single
univariate polynomial $E(T) \in \Q[T]$ of degree four, generating
the same field as $X^4+X^3-1$.

\begin{remark}
The eliminant $E(T)$ is the minimal polynomial of $\tr\rho(\mu)$,
which is a different generator of the same field from the one used in
earlier drafts. Both generate $K_{283}$; the fields are isomorphic and
the discriminant and signature agree. This discrepancy in the
historical record is resolved here by naming the generator explicitly.
\end{remark}

\subsection{Gate F2: certified geometric root}

Certified hyperbolicity at $300$ bits
(\texttt{verify\_hyperbolicity(holonomy=True)}) supplies a rigorous
complex interval containing $\tr\rho_{\mathrm{geom}}(\mu)$. Interval
Newton contraction on $E$ proves that the corresponding box contains
exactly one root of $E$, and that this root is separated from all
others. Hence the geometric character is an exact point of the
eliminated locus.

\subsection{Gate F3: upper inclusion}

The filled character algebra
\begin{equation}
A_{\mathrm{fill}} = \Q[x,y,z]/I_{\mathrm{fill}}
\end{equation}
on the geometric branch has $\dim_\Q A_{\mathrm{fill}} = 4$, and
$I_{\mathrm{fill}}$ is radical and prime. Thus $A_{\mathrm{fill}}$ is
itself a quartic number field, isomorphic to $K_{283}$; no nilpotent
reduction is required. Every word trace of a two-generator $\SLtwo$
representation is a Fricke polynomial in $x,y,z$, so every squared
word trace lies in $K_{283}$, giving
\begin{equation}
\kinv(M) \subseteq K_{283}.
\end{equation}

\subsection{Gate F4: lower inclusion}

For $A = \rho_{\mathrm{geom}}(a)$, Cayley--Hamilton gives
$\tr(A^2) = x^2 - 2$, so $x^2 - 2 \in \kinv(M)$. Exact elimination of
$x$ between the quartic and $S - x^2 = 0$ shows that the minimal
polynomial of $x^2$ has degree four. Hence
$[\Q(x^2):\Q] = 4$, and since $\Q(x^2) \subseteq K_{283}$ with both of
degree four,
\begin{equation}
\Q(x^2) = K_{283} \subseteq \kinv(M).
\end{equation}
Combining the two inclusions proves Theorem~\ref{thm:itf}. \qed

\section{Presentation--geometry bridge}
\label{sec:bridge}

Software-returned filled presentations may replace the literal slope
word by a shortened relator. Any identification of algebraic
presentation coordinates with certified holonomy therefore requires a
bridge.

\begin{proposition}
\label{prop:bridge}
Let $s = \mu^{-2}\lambda^{3}$ be the literal slope word and let $q$ be
the shortened filling relator. Then $q = s^{-1}$ in
$\langle a,b \mid r\rangle$, and consequently
\begin{equation}
\langle\langle r, q \rangle\rangle
= \langle\langle r, s \rangle\rangle
\end{equation}
in the free group on $a,b$. The shortened filled presentation and the
peripheral-slope presentation therefore define the same marked
quotient.
\end{proposition}

The proof is a finite Dehn reduction: explicit rewriting transcripts
reduce both $qs$ and $sq$ to the empty word using cyclic conjugates of
$r$ only. No Knuth--Bendix completion and no finitely-presented-group
equality oracle is invoked.

Certified interval evaluation of the holonomy additionally confirms
that both the cusped relator and the filling relation map to $+I$,
with $-I$ excluded, fixing the $\SLtwo$ lift.

\section{Character-variety structure}
\label{sec:atlas}

\subsection{The target-free atlas}

Before any comparison with physical data, we computed a preregistered
structural atlas of $m003$ and thirteen of its Dehn fillings. The
protocol fixes: the word universe (all freely and cyclically reduced
words in $a,A,b,B$ through length six, quotiented by cyclic rotation
and inversion, proper powers excluded --- $99$ classes); the
generator basis, with a hard guard aborting on drift under filling;
$300$-bit certified hyperbolicity for every filling; and the observable
set. It computes no fitness score and contains no mixing-matrix target.

The atlas recorded four squared-trace degeneracies. Three recur at
every filling and at the cusp. The fourth occurs only at $(-2,3)$.

\subsection{Universal identities}

Let $X = X_{\SLtwo}(m003)$ in the Fricke coordinates above, let $X_0$
be the one-dimensional irreducible component containing the discrete
faithful character, and write $\tau_w = \tr\rho(w)$,
$\sigma_w = \tau_w^2$.

\begin{proposition}[The geometric component is $\mathbb{G}_m$]
\label{prop:x0gm}
Exact primary decomposition of the bare relator ideal
$\langle \rho(r) - I\rangle$ (imposed over $u^2+zu+1=0$, then
eliminating $u$) has three components: two zero-dimensional, with
residue field $\Q[y]/(y^2+y-1)$, and the one-dimensional
\begin{equation}
X_0 = V(P_0), \qquad
P_0 = \langle\, xz + 1,\ \ x^2 + y - 1\,\rangle .
\end{equation}
The further generator $yz - x - z = -x(xz+1) + z(x^2+y-1)$ is
redundant, so $P_0$ is a codimension-two complete intersection. It is
prime, and since $xz = -1$ makes $x$ a unit,
\begin{equation}
\Q[X_0] = \Q[x,y,z]/P_0 \;\xrightarrow{\ \sim\ }\; \Q[x,x^{-1}],
\qquad y \longmapsto 1 - x^2, \quad z \longmapsto -x^{-1}.
\end{equation}
Hence $X_0 \cong \mathbb{G}_m$, uniformized by $t = x = \tr\rho(a)$
through the ring map
\begin{equation}
\Phi : \Q[x,y,z] \twoheadrightarrow \Q[t,t^{-1}], \qquad
(x,y,z) \longmapsto (t,\ 1-t^2,\ -t^{-1}), \qquad
\ker\Phi = P_0 .
\end{equation}
\end{proposition}

\begin{theorem}[Universal identities]
\label{thm:universal}
Put $D_1 = \tau_A - \tau_{ABBB}$, $D_2 = \tau_{AB} - \tau_{ABB}$,
$D_3 = \tau_{AAb} - \tau_{AABB}$. In the Fricke coordinates these
satisfy the polynomial identities
\begin{equation}
D_2 = -(yz - x - z) = x(xz+1) - z(x^2+y-1),
\end{equation}
\begin{equation}
D_1 = (y+1)\,D_2, \qquad
D_3 = (y+1)\bigl[(x^2+y-1) - (xz+1)\bigr],
\end{equation}
in $\Z[x,y,z]$. Consequently $D_1, D_2, D_3 \in P_0$, and
\begin{equation}
\tau_A = \tau_{ABBB}, \qquad
\tau_{AB} = \tau_{ABB}, \qquad
\tau_{AAb} = \tau_{AABB}
\end{equation}
identically on $X_0$; a fortiori $\sigma_A = \sigma_{ABBB}$,
$\sigma_{AB} = \sigma_{ABB}$, $\sigma_{AAb} = \sigma_{AABB}$. The
identities are thus not an unexplained degeneracy of the geometric
character but explicit membership in the defining ideal of $X_0$:
$D_2$ is (up to sign) a defining equation of the component, and
$D_1, D_3$ are $(\tau_B+1)$ times a combination of the two generators
of $P_0$.
\end{theorem}

\begin{proof}
Expanding the Fricke trace polynomials
$\tau_A = x$, $\tau_{ABBB} = -xy + y^2 z - z$,
$\tau_{AB} = z$, $\tau_{ABB} = -x + yz$,
$\tau_{AAb} = x^2 y - xz - y$,
$\tau_{AABB} = -x^2 + xyz - y^2 + 2$
(each reduced modulo $u^2 + zu + 1$) and multiplying out the stated
right-hand sides verifies the three identities in $\Z[x,y,z]$. Each
$D_i$ is then an explicit $\Q[x,y,z]$-combination of the two
generators of $P_0$, hence lies in $P_0$ and vanishes on $X_0$.
Independently, $\Phi$ from Proposition~\ref{prop:x0gm} carries each
pair to a single Laurent function,
\begin{equation}
\Phi(\tau_A) = \Phi(\tau_{ABBB}) = t, \quad
\Phi(\tau_{AB}) = \Phi(\tau_{ABB}) = -t^{-1}, \quad
\Phi(\tau_{AAb}) = \Phi(\tau_{AABB}) = 2t^2 - t^4,
\end{equation}
a one-variable check of the same fact.
\end{proof}

\begin{proposition}[Ambient decomposition and Riley characterization]
\label{prop:jd}
Let $J_D = \langle D_1, D_2, D_3\rangle
= \langle\, yz - x - z,\ \ (y+1)(x^2 + y - xz - 2)\,\rangle$. In the
ambient Fricke $3$-space,
\begin{equation}
\sqrt{J_D}
= P_0 \,\cap\, \langle y+1,\ x+2z\rangle
      \,\cap\, \langle y-2,\ x-z\rangle ,
\end{equation}
an intersection of three primes; that is, $V(J_D) = X_0 \cup L \cup M$
with $L = V(y+1,\,x+2z)$ and $M = V(y-2,\,x-z)$ two affine lines, each
meeting $X_0$ in a finite set and neither contained in it. Restricted
to the relator variety,
\begin{equation}
V\bigl(\langle \rho(r) - I\rangle\bigr) \,\cap\, V(J_D) \;=\; X_0
\end{equation}
set-theoretically: neither $L$ nor $M$ lies on the relator variety,
and the two zero-dimensional relator components satisfy no equation of
$J_D$. The three universal identities therefore characterize the
geometric component inside $X_{\SLtwo}(m003)$.
\end{proposition}

\begin{proof}
From $D_2 = -(yz-x-z)$ and $D_1 = (y+1)D_2$ one has
$J_D = \langle yz-x-z,\ (y+1)(x^2+y-xz-2)\rangle$ (a Gr\"obner basis
confirms the two presentations coincide). With
$g_1 = yz-x-z = -x\,g_2 + z\,g_3$, $g_2 = xz+1$, $g_3 = x^2+y-1$, the
system $g_1 = 0$, $(y+1)(g_3-g_2) = 0$ splits into $y=-1$ (giving $L$)
and $g_3 = g_2$; on $g_1=0$ the latter forces $(z-x)g_2 = 0$, i.e.\
$g_2 = 0$ (giving $X_0$) or $z = x$, which with $g_3-g_2 = y-2 = 0$
gives $M$. The quotients $\Q[x,x^{-1}]$ and $\Q[t]$ are domains, so
$P_0$, $\langle y+1,x+2z\rangle$, $\langle y-2,x-z\rangle$ are prime
and the radical equality follows from an ideal-intersection
computation. Finally $y^2+y-1 \neq 0$ at $y = -1$ and $y = 2$, so $L$
and $M$ meet the relator variety only along $L \cap X_0$,
$M \cap X_0$, and the zero-dimensional relator components do not
satisfy $P_0$; hence the displayed set equality. At the level of the
certified radical decomposition
$\bigl(\bigcap_i C_i\bigr) + J_D = P_0$, with $C_i$ the three relator
components; no scheme-theoretic claim is made for the unreduced ideal
$\langle \rho(r) - I\rangle$.
\end{proof}

\begin{corollary}
Every character obtained by specialization along $X_0$ --- including
the complete cusped character and the geometric characters of all
Dehn fillings on the canonical branch --- satisfies these identities.
Their recurrence across the atlas is one algebraic fact, not fourteen
coincidences.
\end{corollary}

\subsection{The collision ideal and the line $N$}
\label{sec:collision}

The identities of Theorem~\ref{thm:universal} are instances of a general
phenomenon: distinct words whose traces agree on $X_0$. This subsection
determines \emph{all} such coincidences, for words of every length. To
avoid a clash with the manifold $M$ we write $N = V(x-z,\ y-2)$ for the
line called $M$ in Proposition~\ref{prop:jd}; its ideal is
$I(N)=\langle x-z,\ y-2\rangle$. Every point of $N$ is the character of a
representation with $b\mapsto I$, i.e.\ one factoring through the
abelianization $F_2\twoheadrightarrow\Z\langle a\rangle$; we call $N$ the
\emph{abelian (or cyclic-character) specialization line} accordingly ---
more informative than merely ``reducible,'' since $\kappa$ vanishing on
$N$ (Remark~\ref{rem:structure} below) only records reducibility, not
this specific cyclic factorization, which is what the proof actually uses.

For a word $w$ in $a^{\pm1},b^{\pm1}$ let $\tau_w\in\Z[x,y,z]$ be the
Fricke trace polynomial, so that $\tr\rho(w)=\tau_w(x,y,z)$ for every
$\rho:F_2\to\SLtwo$ with $x=\tr\rho(a)$, $y=\tr\rho(b)$, $z=\tr\rho(ab)$;
it is unique because $\rho\mapsto(x,y,z)$ is onto $\mathbb{C}^3$. Let
$e_a(w)\in\Z$ be the exponent sum of $a$ in $w$, and define the
polynomials $S_n\in\Z[x]$ by $S_0=2$, $S_1=x$, $S_{n+1}=xS_n-S_{n-1}$.
For words $w,v$ put $\Delta_{w,v}=\tau_w-\tau_v$, and define the
\emph{collision ideal}
\begin{equation}
J_\infty=\bigl\langle\,\Delta_{w,v}\ :\ \Phi(\tau_w)=\Phi(\tau_v)\,\bigr\rangle
\subseteq\Q[x,y,z],
\end{equation}
where $\Phi$ is the uniformization of Proposition~\ref{prop:x0gm}; by that
proposition $\Phi(\tau_w)=\Phi(\tau_v)$ holds exactly when
$\Delta_{w,v}\in P_0$, and so $J_\infty\subseteq P_0$. Let $J_L$ be the
ideal generated by the same differences with $|w|,|v|\le L$.

\begin{lemma}[Specialization to $N$]
\label{lem:nexp}
For every word $w$, $\ \tau_w(x,2,x)=S_{|e_a(w)|}(x)$ in $\Z[x]$.
Consequently $\Delta_{w,v}\in I(N)$ if and only if $|e_a(w)|=|e_a(v)|$.
\end{lemma}

\begin{proof}
For $c\in\mathbb{C}$ let $\rho_c(a)=\begin{pmatrix}c&-1\\1&0\end{pmatrix}$
and $\rho_c(b)=I$. Then $(x,y,z)=(c,2,c)$, and
$\tr\rho_c(w)=\tr\bigl(\rho_c(a)^{e_a(w)}\bigr)$; the sign of the exponent
is immaterial, since $\tr A^{-1}=\tr A$. The Cayley--Hamilton relation
$A^{k+1}=cA^k-A^{k-1}$ then gives $\tr A^k=S_{|k|}(c)$ for all
$k\in\Z$ (equivalently $S_{-k}=S_k$, as $\lambda^k+\lambda^{-k}$ is
symmetric). Thus $\tau_w(c,2,c)=S_{|e_a(w)|}(c)$ for every
$c\in\mathbb{C}$. As $c$ runs over $\mathbb{C}$ the point $(c,2,c)$ runs
over \emph{all} of $N$, so the two polynomials in $x$
agree at every $c$ and hence identically; no density argument is needed.
Since $S_n$ is monic of degree $n$ for
$n\ge1$ and $S_0=2$, the polynomials $S_n$ are pairwise distinct, so
$\Delta_{w,v}(x,2,x)=S_{|e_a(w)|}-S_{|e_a(v)|}$ vanishes iff
$|e_a(w)|=|e_a(v)|$; and $I(N)$ is the ideal of polynomials vanishing on
$N$, since $\langle x-z,y-2\rangle$ is prime.
\end{proof}

\begin{proposition}[Intersection-point criterion]
\label{prop:ipc}
Let $p=(c,2,c)\in N$ with $c=\lambda+\lambda^{-1}$. Then
$n\mapsto S_n(c)$ is injective on $n\ge0$ if and only if $\lambda$ is not
a root of unity, i.e.\ if and only if $c\notin\{2\cos 2\pi q:\ q\in\Q\}$.
In that case, if $w,v$ are words with $\tau_w(p)=\tau_v(p)$ --- in
particular if $\tau_w=\tau_v$ on any set containing $p$, with no
hypothesis on that set --- then $|e_a(w)|=|e_a(v)|$ and
$\Delta_{w,v}\in I(N)$.
\end{proposition}

\begin{proof}
Since $S_k(c)=\lambda^k+\lambda^{-k}$ (induction on $k$), for $n,m\ge0$
\begin{equation}
\lambda^{n}\bigl(S_n(c)-S_m(c)\bigr)
=\lambda^{2n}+1-\lambda^{n+m}-\lambda^{n-m}
=\bigl(\lambda^{n-m}-1\bigr)\bigl(\lambda^{n+m}-1\bigr).
\end{equation}
If $n\neq m$ then $n-m\neq0$ and $n+m>0$, so the right side vanishes
only if $\lambda$ is a root of unity; conversely if
$\lambda^N=1$ with $N\ge1$, then $S_{N}(c)=S_0(c)$, so injectivity
fails. This proves the equivalence. Now let $\lambda$ be as stated and
$\tau_w(p)=\tau_v(p)$. By Lemma~\ref{lem:nexp} the two values are
$S_n(c)$ and $S_m(c)$ with $n=|e_a(w)|$, $m=|e_a(v)|$, so injectivity
gives $n=m$. \emph{This step turns a single-point equality into an
equality of exponent sums, which are combinatorial invariants of the
words}; the second application of Lemma~\ref{lem:nexp} then upgrades
it to the polynomial identity $\Delta_{w,v}(x,2,x)=S_n-S_m\equiv0$ along
all of $N$, i.e.\ $\Delta_{w,v}\in I(N)$.
\end{proof}

\begin{theorem}[Collision ideal of $X_0(m003)$]
\label{thm:collision}
For all words $w,v$: if $\tau_w=\tau_v$ on $X_0$, then
$|e_a(w)|=|e_a(v)|$ and $\Delta_{w,v}\in I(N)$. Consequently, with
$h=x^2-xz+y-2$ and $q=yz-x-z$,
\begin{equation}
J_\infty \;=\; J_5 \;=\; P_0\cap I(N) \;=\; \langle\, h,\ q\,\rangle .
\end{equation}
Moreover $|\Phi(\tau_w)(i)|=L_{|e_a(w)|}$, where $L_0=2,\ L_1=1,\ L_2=3,\
L_3=4,\dots$ are the Lucas numbers.
\end{theorem}

\begin{proof}
\emph{Intersection point.} Imposing $x=z$, $y=2$ on $P_0$ gives
$x^2=-1$; so $p=(i,2,i)\in X_0\cap N$ (indeed $xz+1=0$ and
$x^2+y-1=0$ there). Here $c=i$ and $\lambda^2-i\lambda+1=0$ gives
$\lambda=i(1\pm\sqrt5)/2$, of modulus $\varphi=(1+\sqrt5)/2\neq1$ or
$1/\varphi\neq1$; a root of unity has modulus $1$, so $\lambda$ is not a
root of unity. If $\tau_w=\tau_v$ on $X_0$ then $\tau_w(p)=\tau_v(p)$, and
Proposition~\ref{prop:ipc} gives the first assertion. (Explicitly,
$S_n(i)=(i\varphi)^n+(-i/\varphi)^n=i^nL_n$, so $|S_n(i)|=L_n$; the
$L_n$ are distinct because $L_0=2$, $L_1=1$ and $L_n$ is strictly
increasing from $n=1$ with $L_n\ge3$ for $n\ge2$. Since
$\Phi(\tau_w)(i)=\tau_w(p)$, this is the final assertion.)

\emph{The ideal $K=P_0\cap I(N)$.} Write $g_2=xz+1$, $g_3=x^2+y-1$.
Then $h=g_3-g_2$ and $q=(z-x)g_2+zh$, and both vanish on $N$, so
$\langle h,q\rangle\subseteq K$. Conversely $P_0=\langle g_2,h\rangle$; if
$ag_2+bh\in I(N)$, restricting to $N$ (where $h=0$ and $g_2=x^2+1$)
gives $a|_N=0$, so $a\in I(N)$. Hence
$K=\langle h\rangle+I(N)g_2$. Now $(x-z)g_2=zh-q$ and
$(y-2)g_2=hg_2-x(x-z)g_2$ (using $y-2=h-x(x-z)$), both in
$\langle h,q\rangle$. So $K=\langle h,q\rangle$.

\emph{Upper bound.} Each generator $\Delta_{w,v}$ of $J_\infty$ lies in
$P_0$ and, by the first assertion, in $I(N)$; hence $J_\infty\subseteq K$.

\emph{Lower bound.} The word pairs $(AB,ABB)$, $(AAb,AABB)$, $(AAbb,AABBB)$
have lengths $\le3,4,5$ and satisfy, in $\Z[x,y,z]$,
\begin{equation}
\Delta_{AB,ABB}=-q,\qquad
\Delta_{AAb,AABB}=(y+1)h,\qquad
\Delta_{AAbb,AABBB}=(y^2+y-1)h,
\end{equation}
(with $\tau_{AAbb}=x^2y^2-x^2-xyz-y^2+2$ and
$\tau_{AABBB}=-x^2y+xy^2z-xz-y^3+3y$; the first two are $D_2,D_3$ of
Theorem~\ref{thm:universal}). Each right-hand side lies in $K\subseteq
P_0$, so each pair collides on $X_0$. Since
$y(y+1)-(y^2+y-1)=1$,
\begin{equation}
h \;=\; y\,\Delta_{AAb,AABB}-\Delta_{AAbb,AABBB}\ \in J_5,
\end{equation}
and $q\in J_5$, so $K=\langle h,q\rangle\subseteq J_5\subseteq J_\infty$.
Combining, $J_\infty=J_5=K$.
\end{proof}

\begin{remark}[Structure of the argument; what $N$ is]
\label{rem:structure}
Two independent jobs are being done. The intersection-point argument
(Proposition~\ref{prop:ipc}) is unconditional and gives
$J_\infty\subseteq P_0\cap I(N)$ for words of every length; a finite
certificate gives $P_0\cap I(N)\subseteq J_5$. Together,
$J_\infty=J_5$. This is a different kind of statement from
``the census through length $10$ stabilizes at $J_5$.'' Also, the
reducibility discriminant $\kappa=x^2+y^2+z^2-xyz-4$ vanishes
identically on $N$ (substitute $z=x$, $y=2$), so $N$ is a line of
reducible characters --- those of representations with $b\mapsto I$ ---
and every point of $X_0\cap N$ is a reducible character lying on the
component. By contrast $\kappa\not\equiv0$ on $X_0$.
\end{remark}

\begin{remark}[Independence from computation]
The proof above uses no enumeration of words. A machine census of all
$4692$ word classes of length $\le10$ (Section~\ref{sec:repro}) serves
only as an independent regression check: it confirms
$\tau_w(i,2,i)=S_{|e_a(w)|}(i)=i^{|e_a|}L_{|e_a|}$ for every word, and
that all $1316$ genuine collision differences at length $\le10$ (from
$1087$ collision groups, $546$ of which are purely ambient, i.e.\
$\tau_w=\tau_v$ already in $\Z[x,y,z]$) lie in $K$, with $J_L=K$ from
$L=5$ on. None of this is logically required by
Theorem~\ref{thm:collision}.
\end{remark}

\begin{remark}[The component of $m006$; other components]
The same argument applies to the geometric component of $m006$,
$X_0=V(x-z,\ yz^2-y-1)$ with $\ker\Phi=\langle x-z,\,yz^2-y-1\rangle$
(checked by elimination). Imposing $x=z$, $y=2$ gives $x^2=3/2$, so
$p=(r,2,r)$ with $r=\sqrt{3/2}\in X_0\cap N$. Here $r$ is not an
algebraic integer ($2r^2-3=0$ is its minimal polynomial), whereas
$\lambda+\lambda^{-1}$ is an algebraic integer whenever $\lambda$ is a
root of unity; so Proposition~\ref{prop:ipc} applies, and every
collision on this $X_0$ has $|e_a(w)|=|e_a(v)|$. Hence its collision
ideal is contained in $P_{006}\cap I(N)$; a Gr\"obner computation gives
$P_{006}\cap I(N)=\langle x-z,\ (y-2)(yz^2-y-1)\rangle=J_4$ (the pairs
$(A,AB)$ and $(AAb,AAbb)$ give $x-z$ and, modulo it,
$-(y-2)(yz^2-y-1)$), so $J_\infty=J_4$ there as well. For any other
component, the criterion needs only a point $(c,2,c)$ on it with
$c=\lambda+\lambda^{-1}$, $\lambda$ not a root of unity; whether such a
point exists must be checked component by component, and we have not
done so beyond $m003$ and $m006$.
\end{remark}

\begin{remark}[Coordinate-free form]
Nothing singles out $a$. For integers $\alpha,\beta$ and
$\ell(w)=\alpha e_a(w)+\beta e_b(w)$, put $\rho(a)=A^{\alpha}$,
$\rho(b)=A^{\beta}$ with $A$ as in the proof of
Lemma~\ref{lem:nexp}. Then $\tr\rho(w)=S_{|\ell(w)|}(c)$, and the
character is $(S_{|\alpha|}(c),S_{|\beta|}(c),S_{|\alpha+\beta|}(c))$,
a rational curve that is the line $N$ for $(\alpha,\beta)=(1,0)$.
The proof of Proposition~\ref{prop:ipc} applies verbatim: if a set $X$
contains such a character with $c=\lambda+\lambda^{-1}$, $\lambda$ not a
root of unity, then $\tau_w=\tau_v$ on $X$ forces $|\ell(w)|=|\ell(v)|$.
(The identity $\tau_w(S_{|\alpha|},S_{|\beta|},S_{|\alpha+\beta|})=S_{|\ell(w)|}$
was also checked, as a regression test, for $(\alpha,\beta)\in\{(1,0),(0,1),
(1,1),(2,3),(1,-2)\}$ over all $4692$ word classes of length $\le10$.)
We use only $\ell=e_a$.
\end{remark}

\begin{remark}[Fox-calculus reading of $X_0\cap N$]
For the relators $r=abAAbabbb$ ($m003$) and $r=ababbAAbb$ ($m006$) one
has $e_a(r)=0$, and for the abelian character $\psi:a\mapsto t$, $b\mapsto1$
the Fox derivatives satisfy $\partial r/\partial a\,(\psi)=0$ and
\begin{equation}
\frac{\partial r}{\partial b}(\psi)=t^{-1}(t^2+3t+1)\ \ (m003),\qquad
\frac{\partial r}{\partial b}(\psi)=2t^2+t+2\ \ (m006).
\end{equation}
Writing $t=\lambda^2$ and $c=\lambda+\lambda^{-1}$, so that
$c^2=t+2+t^{-1}$, the zeros give $c^2=-1$ and $c^2=3/2$ respectively, which
are exactly the abscissae of the points of $X_0\cap N$ found above by
elimination. For $m003$ the roots $t=-\varphi^{\pm2}$ ($\varphi$ the golden
ratio) are not roots of unity because they are real of modulus
$\ne1$; for $m006$ the polynomial $2t^2+t+2$ is primitive and not monic, so its
roots are not algebraic integers. This is an observation about two cases: the
agreement is what the classical necessary condition for a reducible
character to lie on an irreducible component (vanishing of the Fox row at
$\psi$, i.e.\ $H^1(\pi;\mathbb{C}_\psi)\ne0$) would predict, but no
general statement is proved or used here.

Practically, this gives a route to the hypothesis of
Proposition~\ref{prop:ipc} that does not require first computing the
character variety: from a two-generator presentation alone, compute
$\partial r/\partial b$ at $\psi:a\mapsto t,\,b\mapsto1$ (after checking
$e_a(r)=0$, without which no reducible character can lie on the geometric
component at all), solve for its roots $t$, form $c^2=t+2+t^{-1}$, and
test the resulting $c$ for being non-cyclotomic. This locates
\emph{candidates} for $p\in X\cap N$ directly from the presentation; it
does not by itself confirm that a given root corresponds to a point of
the correct component (as opposed to another component of the character
variety, or an extraneous factor of the Fox polynomial), which still
needs an independent check, as done above by elimination for $m003$ and
$m006$.
\end{remark}

\begin{remark}[Related classical results]
That word traces are integral polynomials in $x,y,z$ is due to Horowitz
\cite{Horowitz1972}, who studies when distinct words have the same
character under \emph{all} representations of $F_2$ into $\mathrm{SL}_2$
(not on a subvariety), and shows (Lemma 6.1 of \cite{Horowitz1972}) that
this universal equality forces the multisets of syllable exponents
$|\alpha_i|$, $|\beta_i|$ to agree up to rearrangement. Southcott
\cite{Southcott1979} sharpens this in two ways relevant here. His main
Theorem 4.1 shows trace polynomials fall into $2^n$ classes by exponent
parity mod $2$ (a different, coarser statement than ours). More directly
relevant: immediately before his Theorem 6.8, Southcott records that
universal trace equality $\tau_u=\tau_v$ forces
$|\sum_i\alpha_i|=|\sum_i\gamma_i|$ --- exactly the quantity $|e_a|$ used
here, not merely its parity --- \emph{by the identical mechanism}: setting
one generator to $I$ and evaluating the other via a recursion
($C_0=2,C_1=z,C_m=zC_{m-1}-C_{m-2}$) that is literally our $S_n$. So the
$b\mapsto I$ device and the recovery of $|e_a|$ (not just its parity) from
it are Horowitz/Southcott's, not new here; an earlier draft of this remark
understated this. Ginzburg--Rudnick \cite{GinzburgRudnick1998} give an
independent, arguably cleaner, derivation of the same fact: using the
reducible-representation transversal $A=\mathrm{diag}(a,a^{-1})$,
$B=Z\,\mathrm{diag}(b,b^{-1})Z^{-1}$ (their $x=0$ is the reducible locus),
they compute the trace polynomial's constant term (Proposition 3.2) as
$c_0=a^R b^M+a^{-R}b^{-M}$ with $R=\sum r_j$, $M=\sum m_j$ the exact
signed syllable-exponent sums (our $e_a,e_b$); comparing this Laurent
monomial identity under universal equality forces $(R,M)=\pm(R',M')$
(their Corollary 3.1) directly from algebraic independence of $a,b$ ---
again no injectivity argument, since $a,b$ range freely. Their main
theorem goes further, identifying the complete symmetry group (cyclic
rotation and reversal) realizing trace equality for ``non-singular''
exponent tuples. Wang \cite{Wang2007} studies the same universal-equality
question from a different angle (an explicit equivalence relation on
words via inverse/shift/mirror operators that is sufficient but, by an
explicit counterexample, not necessary for trace equality); this paper
does not use reducible representations and is not more directly relevant
than the others. All four papers work entirely within the free group's
own representation variety; none restricts to a subvariety cut out by a
group relator. What is not in any of them, and is the actual content
of Proposition~\ref{prop:ipc} and Theorem~\ref{thm:collision}: their
hypothesis is universal (equality for \emph{every} representation of the
free group), so the specializing parameter $z=\tau_a(\rho)$ is free as
$\rho$ ranges over all of $\mathrm{SL}_2$, and $S_n(z)=S_m(z)$ as
polynomials in the indeterminate $z$ forces $n=m$ immediately, from the
degree alone --- no injectivity argument is needed. Our hypothesis is
equality restricted to a proper subvariety $X_0$ cut out by a
$3$-manifold's relator; the specialization line $N$ meets $X_0$ only at
isolated points, so the only equality available is the \emph{numerical}
coincidence $S_n(c)=S_m(c)$ at one fixed $c\in\mathbb C$, not a polynomial
identity in a free variable, and matching indices requires the
non-cyclotomic criterion of Proposition~\ref{prop:ipc}, which has no
counterpart in the universal setting. The finite-stabilization statement
$J_\infty=J_5=P_0\cap I(N)$ likewise has no analogue in
\cite{Horowitz1972,Southcott1979,GinzburgRudnick1998,Wang2007}, none of
which considers a relator-cut subvariety or its ideal. We are not aware
of this restricted-component intersection criterion, or of the resulting
collision-ideal closure, elsewhere in the trace-equivalence literature ---
short of asserting it is new, since one further source in this chain,
Thompson's construction of a free-group embedding realizing universal
tr-equivalence as literal conjugacy \cite{Thompson1989}, was accessible to
us only through its introduction (paywalled beyond that point) and so has
not been fully checked.
\end{remark}

\subsection{Characterization of the $(-2,3)$ locus}

Let $B = b^{-1}$ and $Abb = a^{-1}b^2$. In Fricke coordinates
\begin{equation}
\tr\rho(B) = y, \qquad
\tr\rho(Abb) = xy^2 - yz - x,
\end{equation}
the second by Cayley--Hamilton and the standard trace identities.
Set $C_\pm = \tr\rho(Abb) \mp \tr\rho(B)$.

\begin{theorem}[Signed characterization]
\label{thm:signed}
$C_+$ is nonzero on every irreducible component of $X$. Exact primary
decomposition gives
\begin{equation}
V\bigl(I_X + \langle C_+ \rangle\bigr)
= X^{+}_{-2,3} \cup X_{\mathrm{red}},
\end{equation}
where $X_{\mathrm{red}} = V(y+z+1,\; x+z+1,\; z^2+z-1)$ consists
entirely of reducible characters. Hence
\begin{equation}
X^{\mathrm{irr}} \cap V(C_+) = \bigl(X^{+}_{-2,3}\bigr)^{\mathrm{irr}} .
\end{equation}
\end{theorem}

\begin{theorem}[Sign-twisted branch]
\label{thm:twist}
Let $\varepsilon: \pi_1(m003) \to \{\pm 1\}$ be the sign character
inducing $(x,y,z) \mapsto (-x, y, -z)$. Then $\varepsilon$ preserves
the cusped relator but $\varepsilon(s) = -1$ on the slope word. The
nonreducible component of the $C_-$ branch is exactly
$\varepsilon(X^{+}_{-2,3})$, verified by exact ideal equality.
\end{theorem}

\begin{corollary}[Projective characterization]
\label{cor:projective}
On irreducible projective characters,
\begin{equation}
\boxed{\;\tr^2\rho(B) = \tr^2\rho(Abb)
\iff \text{the character lies on the } (-2,3) \text{ filling locus}\;}
\end{equation}
\end{corollary}

\begin{remark}
The minus branch is not a second $\SLtwo$ representation of
$\pi_1(M)$: its filling word maps to $-I$ rather than $+I$. It is the
sign-twisted lift on the cusped group of the same $\PSLtwo$
representation. Both branches project to the same filling condition.
\end{remark}

\subsection{Group-theoretic mechanism}

\begin{theorem}[Conjugacy]
\label{thm:conj}
In $\pi_1(m003(-2,3))$, with $g = BaBA$,
\begin{equation}
g\, B\, g^{-1} = Abb^{-1} .
\end{equation}
Consequently every $\SLtwo$ representation of the filled group
satisfies $\tr\rho(B) = \tr\rho(Abb)$.
\end{theorem}

This supplies the mechanism behind the positive branch of
Theorem~\ref{thm:signed}: the trace equality is not a coincidence of
the geometric character but a consequence of an exact word relation in
the filled fundamental group.

\subsection{Summary of the structure}

The two classes of degeneracy found by the atlas arise from different
algebraic mechanisms:
\begin{equation}
\begin{array}{ccc}
\rho(s) = +I & \Longrightarrow & \tr\rho(B) = \tr\rho(Abb) \\[1mm]
\downarrow \varepsilon & & \downarrow \\[1mm]
(\varepsilon\rho)(s) = -I & \Longrightarrow & \tr\rho(B) = -\tr\rho(Abb)
\end{array}
\end{equation}
while both satisfy $\tr^2\rho(B) = \tr^2\rho(Abb)$. The universal
identities of Theorem~\ref{thm:universal} lie in $I(X_0)$; the
degeneracy $C_+$ does not, and cuts out a proper sublocus.

\section{A Borel-factor construction for lepton mixing}
\label{sec:borel}

This section is downstream of the mathematics above and is reported
with its limitations stated.

\subsection{Construction}

For a word $w$ with loxodromic holonomy, let $\hat n(w) \in S^2$ be the
axis direction obtained from the matrix logarithm. Given a triple
$(w_1,w_2,w_3)$, form the lower-triangular matrix $L$ with unit
diagonal and entries $\ell_{ij}$ for $i > j$, take the $QR$
factorization $L = QR$, and compare $|Q|$ with the PDG 2024 modulus
matrix $|U_{\mathrm{PMNS}}|$ under the permutation-minimized Frobenius
norm
\begin{equation}
\mathcal{F} = \min_{\pi \in S_3}
\bigl\| |Q|_\pi - |U_{\mathrm{PMNS}}| \bigr\|_F .
\end{equation}

\subsection{Reported value}

With the triple $\{\texttt{aa}, \texttt{aaB}, \texttt{baa}\}$ on $M$,
optimizing $(\ell_{21}, \ell_{31}, \ell_{32})$ by Nelder--Mead gives
\begin{equation}
\mathcal{F}_{\mathrm{PMNS}} = 0.005087,
\end{equation}
at $\ell_{21} = -1.13114$, $\ell_{31} = -1.02174$,
$\ell_{32} = 1.09631$. This value is reproducible in the recorded
environment.

Of the three parameters, only $\ell_{32}$ has an evident closed-form
reading: $\ell_{32} = 1.09631$ against $\tan\theta_{23} = 1.09665$, a
relative difference of $0.03\%$. The other two have no simple
relationship to the mixing angles.

\subsection{Relation to the atlas}

Under the canonical word convention of Section~\ref{sec:atlas},
$\texttt{aaB} \sim \texttt{AAb}$, which is one member of the universal
identity $\sigma_{AAb} = \sigma_{AABB}$ of
Theorem~\ref{thm:universal}. The word $\texttt{aa}$ does not appear as
an atlas class because it is a proper cyclic power.

We report this overlap as a structural observation. It does not
establish that the identity explains the fit, nor that it selects the
triple.

\subsection{Historical provenance of the Borel construction}
\label{sec:provenance}

The construction reported above is not the first one proposed for this
manifold. An earlier draft (March 2026) defined a genuinely
manifold-dependent, one-parameter construction:
\begin{equation}
\ell_{ij} = \lambda\, s_{ij}\, (\hat n_i \cdot \hat n_j),
\end{equation}
with a single scale $\lambda \in [0.1,5.0]$ optimized over a grid of
$50$ values and discrete signs $s_{ij} \in \{\pm1\}$ optimized over $8$
combinations --- unlike the free three-parameter fit of
Section~\ref{sec:borel}, this family has only one continuous degree of
freedom once the word triple and manifold are fixed, and would support
a substantially narrower null distribution.

That draft reported, for the triple $\{\texttt{aa},\texttt{ab},
\texttt{aB}\}$ on $M$, $\lambda=1$ exactly and signs $(+,-,+)$, giving
$(\ell_{21},\ell_{31},\ell_{32}) \approx (0.443,-0.530,0.432)$ and
fitness $\mathcal F \approx 0.019$.

This specific numerical claim is not currently reproducible. We
searched all sixteen boolean \texttt{fundamental\_group\_args}
conventions available in the present SnapPy release, plus the plain
default, computing the live axis dot products for
$\{\texttt{aa},\texttt{ab},\texttt{aB}\}$ under each; none comes within
an order of magnitude of $(0.443,-0.530,0.432)$ (closest:
$(0.604,0.379,-0.320)$, still far from the target in both magnitude
and pattern). The historical generator convention has therefore not
been recovered by this search, and no more specific hypothesis (a
differing SnapPy version, an unrecorded triangulation choice) has been
confirmed. This is consistent with, and does not go beyond, a
contemporaneous note in the project's own history flagging that these
particular words used ``a different generator convention'' without
further diagnosis.

We therefore make no numerical claim from the March 2026 construction.
It remains of interest --- being genuinely manifold-dependent, it
would not be subject to the manifold-independence result that applies
to the free three-parameter fit --- but recovering its exact
generator convention is a separate, unresolved archival problem, not
a mathematical one.

\section{Null calibration}
\label{sec:null}

An earlier version of this work reported $p < 10^{-4}$ for the fitness
above, based on comparing Haar-random targets against a construction
with the optimized parameters held fixed. That test is not matched to
the fitting procedure: since $(\ell_{21}, \ell_{31}, \ell_{32})$ are
optimized when fitting the physical target, a null trial must repeat
that optimization for every random target.

Repeating the calibration correctly, with independent optimization per
target, over $N = 50000$ Haar-random unitary targets gives
\begin{equation}
k = 6625, \qquad \hat p = k/N = 0.1325,
\end{equation}
with binomial standard error $0.0015$. An interim run at $N = 2000$
gave $250/2000 = 0.125$, consistent.

\begin{remark}
The reason is dimensional. The map
$(\ell_{21},\ell_{31},\ell_{32}) \mapsto |Q|$ has full rank $3$ at the
fitted point, with condition number $\approx 4.1$. A three-parameter
family covering a three-dimensional target manifold is expected to
approach random targets at this rate. No further mechanism is needed
to explain the result.
\end{remark}

\begin{corollary}
The claim $p < 10^{-4}$ is withdrawn. The fitness
$\mathcal{F}_{\mathrm{PMNS}} = 0.005087$ remains a reproducible
property of the construction, but under matched-procedure calibration
it does not constitute evidence that the physical target is
statistically exceptional.
\end{corollary}

\section{Open questions}

\begin{enumerate}
\item \textbf{CP phase basis dependence.} A twist-angle formula
$\delta = \pi + \varphi(\texttt{aaB}) + \varphi(\texttt{baa})$ gives
$195.91^\circ$ under one generating set and $123.96^\circ$ under the
verified cusped basis of Section~\ref{sec:homology}. The words
$\texttt{aaB}$ and $\texttt{baa}$ are not well defined without
specifying the generating set. Until this is resolved the formula is
not well posed and no value is asserted.

\item \textbf{Constrained geometric construction.} See
Section~\ref{sec:provenance}: an earlier, genuinely manifold-dependent
one-parameter construction is documented in the historical record but
its exact generator convention has not been recovered despite a
systematic search, so it supports no numerical claim here.

\item \textbf{Physical selection.} Nothing in Sections
\ref{sec:homology}--\ref{sec:atlas} establishes that $M$ must encode
lepton mixing. The exact results describe the manifold; they do not
supply a derivation of the phenomenology.
\end{enumerate}

\section{Reproducibility}
\label{sec:repro}

All computations are exact or interval-certified as indicated.
Scripts and logs are archived with SHA-256 hashes in the project
repository.

\begin{table}[h]
\centering
\caption{Certificate inventory.}
\begin{tabular}{llc}
\toprule
Result & Script & Status \\
\midrule
Surgery law (Thm.~\ref{thm:surgery}) & \texttt{linking\_form.py} & exact \\
ITF gates F1--F2 (Thm.~\ref{thm:itf}) & \texttt{pmns\_itf\_certificate.sage} & exact + certified \\
ITF gates F3--F4 (Thm.~\ref{thm:itf}) & \texttt{pmns\_itf\_certificate.sage} & exact \\
Presentation bridge (Prop.~\ref{prop:bridge}) & \texttt{ckm\_dehn\_relator\_certificate} & exact \\
Target-free atlas & \texttt{m003\_structure\_atlas.sage} & computed \\
Universal identities (Thm.~\ref{thm:universal}) & \texttt{m003\_three\_universal\_identities.sage} & exact \\
Component $\cong \mathbb{G}_m$, ambient decomp.\ (Prop.~\ref{prop:x0gm}, \ref{prop:jd}) & \texttt{m003\_stage4\_x0\_identity\_origin.py} & exact \\
Collision ideal (Lem.~\ref{lem:nexp}, Thm.~\ref{thm:collision}) & \texttt{m003\_m006\_X0\_cap\_N\_point\_argument.py} & exact \\
$J_L=K$, $L\le10$ ($m003$; regression check) & \texttt{m003\_J10\_vs\_intersection.py} & exact \\
$J_L=K$, $L\le10$ ($m006$; regression check) & \texttt{m006\_J10\_vs\_intersection.py} & exact \\
Signed characterization (Thm.~\ref{thm:signed}) & \texttt{m003\_squared\_locus\_and\_conjugacy.sage} & exact \\
Sign twist (Thm.~\ref{thm:twist}) & \texttt{m003\_identify\_Y\_minus.sage} & exact \\
Conjugacy (Thm.~\ref{thm:conj}) & \texttt{m003\_squared\_locus\_and\_conjugacy.sage} & exact \\
Borel fitness (\S\ref{sec:borel}) & \texttt{hfg\_reproduce.py} & computed \\
Null calibration (\S\ref{sec:null}) & \texttt{pmns\_null\_corrected.py} & computed \\
\bottomrule
\end{tabular}
\end{table}

The atlas protocol, including its word universe, generator-basis guard,
precision, and control filling families, was frozen and hashed before
the structural results of Section~\ref{sec:atlas} were examined.

\section*{Acknowledgments}

The author used AI language model assistance (Claude, Anthropic) for
editorial support in preparing this manuscript. All mathematical
content, computations, and scientific claims are the author's own
work, independently verified using SnapPy and SageMath.

\begin{thebibliography}{9}

\bibitem{GMM2009}
D.~Gabai, R.~Meyerhoff, and P.~Milley,
\emph{Minimum volume cusped hyperbolic three-manifolds},
J. Amer. Math. Soc. \textbf{22} (2009), 1157--1215.

\bibitem{MR}
C.~Maclachlan and A.~W.~Reid,
\emph{The Arithmetic of Hyperbolic 3-Manifolds},
Graduate Texts in Mathematics \textbf{219}, Springer, 2003.

\bibitem{SnapPy}
M.~Culler, N.~M.~Dunfield, M.~Goerner, and J.~R.~Weeks,
\emph{SnapPy, a computer program for studying the geometry and
topology of 3-manifolds}, \url{http://snappy.computop.org}.

\bibitem{Sage}
The Sage Developers,
\emph{SageMath, the Sage Mathematics Software System}, 2026,
\url{https://www.sagemath.org}.

\bibitem{PDG}
Particle Data Group,
\emph{Review of Particle Physics},
Prog. Theor. Exp. Phys. (2024).

\bibitem{Horowitz1972}
R.~D.~Horowitz,
\emph{Characters of free groups represented in the two-dimensional
special linear group},
Comm. Pure Appl. Math. \textbf{25} (1972), 635--649.

\bibitem{Southcott1979}
J.~B.~Southcott,
\emph{Trace polynomials of words in special linear groups},
J. Austral. Math. Soc. Ser. A \textbf{28} (1979), 401--412.

\bibitem{GinzburgRudnick1998}
D.~Ginzburg and Z.~Rudnick,
\emph{Stable multiplicities in the length spectrum of Riemann surfaces},
Israel J. Math. \textbf{104} (1998), 129--144.

\bibitem{Wang2007}
G.~Wang,
\emph{A note on trace polynomial},
Tsinghua Sci. Technol. \textbf{12} (2007), no.~4, 479--484.

\bibitem{Thompson1989}
J.~G.~Thompson,
\emph{Fricke, free groups, and $SL_2$},
in Group Theory: Proceedings of the Singapore Group Theory Conference,
de Gruyter, Berlin, 1989, 207--214.

\bibitem{Porti}
J.~Porti,
\emph{Character varieties and knot symmetries},
lecture notes, 2017.

\end{thebibliography}

\end{document}
